\scanpage{44}
\margin{43}
\mspage{1}
\margin{Attention, notations \unsure{de ce qui suit} p.r. à celles des calculs précédents, en $\alpha,\beta,\alpha',\beta',\xi,\eta,\xi'\ldots$}
Considérons une 2-sphère conforme \inserted{$\Sigma$} avec une partie $J$ à trois éléments,
\[
(1)\qquad J\subset\Sigma,\qquad J=(R_i)_{i\in J}
\]
de sorte que le groupe $\unsure{\mathcal G_J|G_J}$ opère sur $\Sigma$ par automorphismes \deleted{orientés} directs, \deleted{et que} \deleted{$\unsure{G'_J|G_J}\times\{1,\tau\}$ le groupe des} \inserted{tandis que le groupe des automorphismes} conformes de la situation $(\Sigma,J)$ est
\[
(2)\qquad\Gamma_{\Sigma,J}=\Gamma\simeq\unsure{\mathcal G_J|G_J}\times\{1,\tau\}
\]
\margin{\uline{NB} L'antipodisme ne fait \uline{pas} partie de $\Gamma$, car il ne stabilise pas $J$}
$\tau$ est l'unique \deleted{anti} autom. $\unsure{+|\ne}1$ \deleted{\unsure{qui est}} qui fixe les $R_i$ \deleted{\unsure{tous}} ($\tau={}$conjugaison complexe).
\margin{donc la \unsure{catégorie des} $(\Sigma,J)$ équiv. à celle des couples $(J,\uline{\bar\omega})$, où $J\in(\mathrm{Ens})_3$ et $\uline{\bar\omega}\in(\mathrm{Ens})_2$ cf plus bas}
Il y a sur $\Sigma$ une unique métrique riemannienne homogène \unsure{invariante|invariant} pour $\unsure{\mathcal G_J|G_J}$, \inserted{à scalaire $\lambda\in\mathbb R^{++}$ près} \deleted{pour laquelle}, qui nous permet d'identifier canoniquement $\Sigma$ à la sphère unité dans un espace vectoriel réel euclidien $E_{\unsure{3}}$. On pose
\[
(3)\qquad\left\{\begin{aligned}
\uline{\bar\omega}&=\Omega(\Sigma)\simeq\Omega(E_{\unsure{3}})\\
\uline{\omega}&=\Omega(J)\simeq\Omega(\Sigma_{\mathbb R})\\
\uline{\varepsilon}&=\Pi_0(\Sigma\setminus\Sigma_{\mathbb R})=\Pi_0(E\setminus\unsure{V|\forall})=\Pi_0((E/V)^*)
\end{aligned}\right.
\]
\margin{ens des deux orientations de $\Sigma$, ou donc de $E$}
\margin{On pourrait écrire $E(J,\uline{\bar\omega})$ $V(J,\uline{\omega})$.}
\margin{On a
\[
\uline{\omega}\wedge\uline{\varepsilon}\wedge\uline{\bar\omega}=1
\]
i.e.
\[
(3\text{ bis})\qquad\left\{\begin{aligned}
\uline{\omega}&\simeq\uline{\varepsilon}\wedge\uline{\bar\omega}\\
\uline{\varepsilon}&\simeq\uline{\bar\omega}\wedge\uline{\omega}\\
\uline{\bar\omega}&\simeq\uline{\omega}\wedge\uline{\varepsilon}
\end{aligned}\right.
\]
par la conv. de \unsure{Stokes entre} orientations d'un hémisphère et celles de \unsure{son|leur} bord, \unsure{plus précis (4) ici}}
$V$ est le \inserted{$\mathbb R$-}sous-espace vectoriel de $E$ \uline{engendré} par les $R_{\unsure{J|i}}$, qui est de dimension 2, et où
\[
(4)\qquad\Sigma_{\mathbb R}=\Sigma\cap V
\]
cercle euclidien \unsure{ds le|dans le} plan euclidien $V$
